\clearpage
\section*{September 9, 2026: Missing companions and Third Avenue's unit schedule}
\textit{Agent-drafted follow-up; production corrections remain pending.}

The October 2025 environmental submissions contain direct evidence of a three-building, 297-unit proposal at Third Avenue. Figure 3, on PDF page 77 of both the 4121 work plan and the revised 4133 work plan, labels 4121, 4133 and 4137 as 15-story buildings with 99 apartments each. The accompanying floor schedule adds to 99: five apartments on the first floor, seven on each of floors 2--11, and six on each of floors 12--15.

This clarifies the preceding entry's unresolved environmental count. The reports' prose says 90, but their attached drawings say 99. The July work plan for 4137 contains an older drawing showing three 13-story, 81-unit buildings, despite prose describing 15 stories and 90 units. These environmental reports mix inconsistent design information. Their undated drawing excerpts support the proposed program in the October submissions; they do not establish a DOB-approved schedule or the exact first-filed choice.

The saved DOB extract still records 98, 97 and 99 in the three I1 rows. Foundation and structural filings report 99 for 4121, and a structural filing reports 99 for 4133. All are records in the July 10, 2026 snapshot. An I1 filing type and an original filing date do not make its other fields an immutable historical snapshot. We should preserve the conflicting fields and distinguish the documented 297-unit design from a first-filed or approved measure.

The seven missed companions have identifiable causes. Four filing lots---290 Bergen, 278 Butler, 139 Beach 29th and 35-11 42nd---are absent from the older map used for post-policy linkage. Third Avenue's lots are present, but a 7.239-metre gap between the mapped lots defeats exact touching. River North's 8 Stuyvesant already appears as a same-owner nearby candidate, without an accepted link. The historical route excludes 3 West before linking because its prefiling lot match fails; the DOB route already links it to 7 Wharf.

Code can expose all these failures and generate review candidates. It should recover lot histories, recognize explicit development-lot references, and keep missing land covariates from silently removing constituents. Common ownership and proximity alone should not automatically establish a parent. A small versioned ledger should retain supported pair decisions and their dated sources, with checks against the three disputed pairs. This is a recommendation, not an adopted expansion of the linkage method.

Estimation remains on hold at Jacob's request. Production membership and units are unchanged. The seven-case trace has seven unique rows and twelve fields, with no missing or repeated companion keys. Remaining address and phase ambiguities, production pair decisions, and a consistent unit measurement rule still prevent certification of the final parent panel.

\textit{Reproduction:} Make in \path{tasks/audits/audit_estimation_parent_links/code}; trace and plan rendering in its output directory. Public source PDFs are retained by \path{tasks/audits/download_parent_review_documents}. The detailed review records the document URLs and conflicting schedules. Branch \path{simplify_empirical_tasks}, working changes after 8458668.
\clearpage
\section*{Third Avenue: attached October 2025 development plan}
\includegraphics[width=\textwidth]{../tasks/audits/audit_estimation_parent_links/output/third_avenue_plan.png}

Source: 4121 Third Avenue Remedial Action Work Plan, October 2025 submission, Figure 3, PDF page 77. Each building is labeled 99 units; the floor schedule is at right. This is a proposed design attached to an environmental submission, not a certified DOB-approved schedule.
